\chapter{可测函数}

\begin{solution}{measurable-rational-thresholds}
若 $\{f>q\}$ 对所有 $q\in\Q$ 可测，则对任意 $a\in\R$，
\[
 \{f>a\}=\bigcup_{q\in\Q,\ q>a}\{f>q\}.
\]
右侧是可数并。半直线判据于是给出 $f$ 可测；其他严格或非严格版本同理。
\end{solution}

\begin{solution}{measurable-vector-map}
有理端点开矩形构成 $\R^n$ 的可数拓扑基。矩形
$R=\prod_j(a_j,b_j)$ 的原像是
$\bigcap_j\{a_j<f_j<b_j\}$，可测。任意开集是基元素的至多可数并，所以
其原像可测；开集生成 Borel $\sigma$-代数，故向量映射可测。
\end{solution}

\begin{solution}{measurable-algebra-operations}
映射 $F=(f,g):X\to\R^2$ 可测。函数
\[
 (u,v)\longmapsto u+v,qquad
 (u,v)\longmapsto uv,qquad
 u\longmapsto |u|,qquad
 (u,v)\longmapsto\max\{u,v\}
\]
都连续，因而 Borel 可测；分别与 $F$ 或其坐标函数复合便得结论。有限最小值
由 $\min(u,v)=-\max(-u,-v)$ 得到。
\end{solution}

\begin{solution}{measurable-limsup-formula}
由“对每个 $N$，尾部中仍有一项大于 $a$”可得
\[
 \{\limsup_nf_n>a\}
 =\bigcup_{r\in\Q,\ r>a}
   \bigcap_{N=1}^{\infty}\bigcup_{n\ge N}\{f_n>r\}.
\]
右侧只含可数并、交和可测水平集，故可测。也可先证明尾上确界可测，再取
递减极限。
\end{solution}

\begin{solution}{measurable-convergence-set}
有限实值情形下，$f_n(x)$ 收敛当且仅当它是 Cauchy 列。因此收敛点集为
\[
 \bigcap_{k=1}^{\infty}\bigcup_{N=1}^{\infty}
 \bigcap_{m,n\ge N}\{|f_m-f_n|<1/k\},
\]
可测。扩充实值情形可用
$\liminf f_n=\limsup f_n$ 表示，并另外允许共同值为无穷。
\end{solution}

\begin{solution}{simple-function-canonical-form}
设简单函数的不同取值为 $a_1,\ldots,a_m$。令
$E_k=\{x:\varphi(x)=a_k\}$，则 $E_k$ 可测、两两不交且覆盖 $X$，并有
$\varphi=\sum_ka_k\mathbf1_{E_k}$。任何按不同系数和不交水平集的表示都必须
把 $a_k$ 的水平集写成对应集合，所以规范形式唯一，允许删去空水平集。
\end{solution}

\begin{solution}{nested-dyadic-simple-approximation}
对 $f\ge0$ 定义
\[
 \varphi_n(x)=2^{-n}
 \min\bigl\{\lfloor2^nf(x)\rfloor,,2^{2n}\bigr\},
\]
约定 $\lfloor2^n\cdot(+\infty)\rfloor=+\infty$。该函数只取有限多个值，且
由可测水平集构成，故为简单可测函数。二进网格逐次细化，同时截断高度由
$2^n$ 增至 $2^{n+1}$，所以
\[
 0\le\varphi_n\le\varphi_{n+1}\le f.
\]
若 $f(x)<\infty$，当 $2^n>f(x)$ 后不再发生截断，且
$0\le f(x)-\varphi_n(x)<2^{-n}$；若 $f(x)=+\infty$，则
$\varphi_n(x)=2^n\to+infty$。因此 $\varphi_n\uparrow f$ 逐点成立。
\end{solution}

\begin{solution}{bounded-simple-uniform-approximation}
若 $|f|\le M$，把 $[-M,M]$ 分成长度不超过 $2^{-n}$ 的有限区间，并把
$f(x)$ 向下取到相应左端点，得到简单可测函数 $\varphi_n$。对所有 $x$，
$|f(x)-\varphi_n(x)|\le2^{-n}$，故一致收敛。
\end{solution}

\begin{solution}{ae-modification-complete-space}
若 $g=f$ 于 $X\setminus N$，则
$\{g>a\}\triangle\{f>a\}\subseteq N$。完备性使该对称差可测，故
$\{g>a\}$ 可测。反例取非完备空间中可测零集 $N$ 的不可测子集 $Z$；零
函数 $f$ 可测，令 $g=\mathbf1_Z$，则 $g=f$ a.e.，但
$\{g>1/2\}=Z$ 不可测。
\end{solution}

\begin{solution}{indicator-measurability}
若 $E$ 可测，指标函数的任意半直线原像只能是
$\varnothing,E,E^c$ 或 $X$，故可测。反之，若 $\mathbf1_E$ 可测，则
$E=\{\mathbf1_E>1/2\}$ 可测。
\end{solution}

\begin{solution}{lusin-simple-functions}
写 $\varphi=\sum_{k=1}^ma_k\mathbf1_{E_k}$，其中 $E_k$ 两两不交且 $E$
有限测度。由内正则性，在每个 $E_k$ 中取紧集 $K_k$，使删去测度总和小于
$\varepsilon$。不同 $K_k$ 是有限个互不相交紧集，彼此距离为正；因此
$\varphi$ 在 $K=\bigcup_kK_k$ 上局部常值，从而连续。
\end{solution}

\begin{solution}{lusin-proof-measure-budget}
先删去 $E\setminus E_M$ 少于 $\varepsilon/4$，再取
$K_0\subset E_M$ 删去少于 $\varepsilon/4$。第 $n$ 个简单函数水平集的
紧化删去少于 $\varepsilon2^{-n-2}$ 即可；其总和至多
$\varepsilon/4$。总损失小于 $3\varepsilon/4<\varepsilon$。每个 $K_n$
是有限个紧集之并，故紧；$K_0\cap\bigcap_nK_n$ 是紧集 $K_0$ 中闭集的交，
仍紧。
\end{solution}

\begin{solution}{lusin-continuous-extension}
$K$ 是 $\R^d$ 的闭子集，$f|_K$ 连续。若先在 Lusin 证明中限制
$|f|\le M$，Tietze 延拓定理给出连续 $F:\R^d\to[-M,M]$，满足
$F|_K=f|_K$。一般 Lusin 定理也已先删去大值部分，故同样适用。
\end{solution}

\begin{solution}{measurable-almost-continuous-not-continuous-ae}
Dirichlet 函数 $f=\mathbf1_\Q$ 在每一点不连续。给定有限测度区间和
$\varepsilon>0$，可取开集 $O\supseteq\Q\cap[a,b]$ 且 $m(O)<\varepsilon$。
在闭集 $K=[a,b]\setminus O$ 上 $f\equiv0$，故连续。于是 Lusin 性质成立，
但不连续点集是整个区间，绝非零测。
\end{solution}

\begin{solution}{measurable-functions-generated-sigma}
记题中右侧为 $\mathcal A_0$。它按定义使每个 $f_i$ 的所有 Borel 集原像
可测，所以全部 $f_i$ 对 $\mathcal A_0$ 可测。若 $\mathcal A$ 是另一使
全部 $f_i$ 可测的 $\sigma$-代数，则每个生成元
$f_i^{-1}(B)$ 都属于 $\mathcal A$，故
$\mathcal A_0\subseteq\mathcal A$。这证明最小性。
\end{solution}
